% Part: applied-modal-logics
% Chapter: epistemic-logic
% Section: properties-accessibility

\documentclass[../../../include/open-logic-section]{subfiles}

\begin{document}

\olfileid[pt]{aml}{el}{acc}

\olsection{Relações de acessibilidade e princípios epistêmicos}

Dado o que já sabemos sobre correspondência de frames nas lógicas
modais normais, podemos examinar a forma assumida pelas !!{formula}s
características sob interpretações epistêmicas. Já dissemos que as
lógicas epistêmicas costumam ser interpretadas em S5. Vejamos, então,
como vários princípios epistêmicos são representados e como
correspondem a diferentes condições sobre frames.

Recorde, da lógica modal normal, que diferentes !!{formula}s modais
caracterizavam diferentes propriedades das relações de acessibilidade.
A tabela destaca algumas que correspondem a princípios epistêmicos
específicos.

\begin{table}[t]
    \begin{tabular}{| p{.48\textwidth} || p{.48\textwidth} |}
      \hline
      {\emph{Se $R$ é \dots}} & {\emph{então \dots é verdadeira em~$\mModel{M}$:}} \\
      \hline \hline
        
      & $\Knows (p \lif q) \lif (\Knows p \lif \Knows q)$ 
      \hfill \newline{(Fechamento)} \\
      \hline
      \emph{reflexiva}: $\forall w Rww$  
      & $\Knows p \lif p$ \hfill \newline{(Veridicidade)} \\
      \hline
      \emph{transitiva}: \newline
      $\forall u \forall v \forall w ((Ruv \land Rvw) \lif Ruw)$ & 
      $\Knows p \lif \Knows \Knows p$ \hfill 
           \newline{(Introspecção positiva)} \\
      \hline 
      \emph{euclidiana}: \newline
      $\forall w \forall u \forall v ((Rwu \land Rwv) \lif Ruv)$ &  
      $\lnot \Knows p \lif \Knows \neg \Knows p$ \hfill 
        \newline{(Introspecção negativa)} \\
      \hline
    \end{tabular}
    \caption{Quatro princípios epistêmicos.}
    \ollabel{tab:four}
  \end{table} 

A veridicidade, correspondente ao axioma $T$, costuma ser tratada como
o menos controverso desses princípios, pois representa a afirmação de
que, se !!a{formula} é conhecida, então deve ser verdadeira. O
fechamento e as introspecções positiva e negativa são muito mais
controversos. 

O fechamento, correspondente ao axioma $K$, representa a ideia de que
o conhecimento de um agente é fechado sob implicação. Isso pode parecer
plausível em alguns casos. Por exemplo, posso saber que, se estou em
Victoria, então estou na ilha de Vancouver. Salvo cenários céticos
inusitados, de fato sei que estou em Victoria, e isso também deveria
sugerir que sei que estou na ilha de Vancouver. Nesse caso, portanto,
o fechamento lógico de meu conhecimento pode parecer relativamente
intuitivo. Por outro lado, nem sempre refletimos sobre as consequências
do que sabemos, e isso pode levar a resultados menos intuitivos em
outros casos.

A introspecção positiva, às vezes conhecida como princípio KK, é
formulada como a afirmação de que, se sei algo, então sei que sei. Ela
é a contraparte epistêmica do axioma 4. Correspondentemente, a
introspecção negativa é formulada como a afirmação de que, se eu
\emph{não} sei algo, então sei que não o sei; ela é a contraparte do
axioma 5. Ambas parecem admitir contraexemplos bastante comuns, nos
quais não tenho certeza se sei ou não algo que, de fato, sei. 


\end{document}
